\section{Introduction}
The intersection of health and labour economics has long captivated researchers and policymakers alike. This paper attempts to answer a fundamental question in health and labour economics: are individuals characterized by fundamentally different health types, and what is the association between any such health types and labour market outcomes? Furthermore, I investigate whether health-related factors can serve as predictors of earning potential. These inquiries extend beyond academic interest, carrying profound implications for individuals, healthcare systems, and economic policy.

Using an unsupervised machine learning approach pioneered in \citet{denardi2024wp} and drawing from the rich dataset of the Understanding Society: UK Household Longitudinal Study (UKHLS), I aim to uncover patterns and relationships that illuminate the complex interplay between health trajectories and economic outcomes. My analysis focuses particularly on chronic illnesses and substantial difficulties caused by health problems or disabilities, seeking to identify distinct health types. Through this investigation, I aspire to contribute to a more nuanced understanding of health-related economic disparities and inform policies that can better address the diverse needs of individuals across different health trajectories.

To address these questions, my study is structured around three key areas of investigation:
\begin{enumerate}
    \item \textbf{Frailty trajectories.} I examine the frailty trajectories associated with different health types. By employing unsupervised clustering techniques on longitudinal health data, I identify distinct patterns in how individuals' health evolves over time, particularly focusing on the progression of chronic illnesses and disabilities. This analysis reveals three primary health types, each with unique implications for long-term health outcomes.
    \item \textbf{Employment dynamics.} I explore the employment dynamics across these health types, considering factors such as earnings and employment status. This section illuminates the association between health trajectories and labour market participation, income levels, and human capital accumulation.
    \item \textbf{Predictive power.} I assess the explanatory power of health types in predicting frailty and, by extension, labour market outcomes. By comparing models with and without health type dummy variables, I demonstrate a substantial improvement in the adjusted $R^2$ value, underscoring the significance of considering heterogeneous health trajectories in labour economic analyses.
\end{enumerate}

Through these three investigations, I aim to provide a comprehensive understanding of how chronic illnesses and health trajectories shape economic outcomes. My findings offer valuable insights that can inform both policy decisions and further research in labour economics.
